% Generated draft: requires mathematical and source review.
\documentclass[11pt]{article}
\usepackage[margin=25mm]{geometry}
\usepackage{fontspec}
\setmainfont{Latin Modern Roman}
\usepackage{amsmath,amssymb,mathtools,unicode-math}
\setmathfont{Latin Modern Math}
\usepackage{enumitem,needspace,xurl,hyperref}
\setlength{\parindent}{0pt}
\setlength{\parskip}{5pt}
\setlength{\emergencystretch}{3em}
\newcommand{\N}{\mathbb{N}}
\newcommand{\Z}{\mathbb{Z}}
\newcommand{\Q}{\mathbb{Q}}
\newcommand{\R}{\mathbb{R}}
\newcommand{\C}{\mathbb{C}}
\newcommand{\F}{\mathbb{F}}
\newcommand{\A}{\mathbb{A}}
\newcommand{\PP}{\mathbb{P}}
\newcommand{\E}{\mathbb{E}}
\newcommand{\eps}{\varepsilon}
\newcommand{\dd}{\,\mathrm{d}}
\begin{document}
\section*{UnsolvedMath ID: 701138}
\textbf{Source problem code:} catalogue entry
\subsection{Definitions and mathematical statement}
Let \(p\) be a prime and let \(A\) be an abelian group of exponent \(p\), written additively: \(A\neq\{0\}\) and \(pa=0\) for every \(a\in A\). No finiteness assumption is imposed on \(A\). Let \(G\leq\operatorname{Aut}(A)\). For \(x,g\in G\), use the conventions
\[
[x,g]=x^{-1}g^{-1}xg,\qquad x^g=g^{-1}xg,
\]
and write
\[
\langle x^G\rangle=\bigl\langle g^{-1}xg\mid g\in G\bigr\rangle
\]
for the normal closure of \(x\) in \(G\). A group is locally finite if each of its finitely generated subgroups is finite.

The supplied record does not define the word ``regular.'' Under its fixed-point-free interpretation, this means
\[
C_A(x):=\{a\in A:x(a)=a\}=\{0\}.
\]
With this interpretation, the question of V. D. Mazurov is the following: for every prime \(p\), every abelian group \(A\) of exponent \(p\), every subgroup \(G\leq\operatorname{Aut}(A)\), and every \(x\in G\) satisfying
\[
\operatorname{ord}(x)=3,\qquad C_A(x)=\{0\},\qquad
\forall g\in G\ \exists n_g\in\mathbb Z_{\geq1}\quad [x,g]^{n_g}=1,
\]
is \(\langle x^G\rangle\) locally finite? Equivalently, must
\[
\left|\langle h_1,\ldots,h_r\rangle\right|<\infty
\quad\text{for every }r\geq1\text{ and all }h_1,\ldots,h_r\in\langle x^G\rangle?
\]
The finite orders of the commutators may depend on \(g\); no common bound is assumed. Confirmation of the intended meaning of ``regular'' is required for this conditional formulation.
\subsection{Short English statement}
Does a fixed-point-free automorphism \(x\) of order \(3\) of an abelian group of prime exponent generate a locally finite normal subgroup of \(G\), provided every commutator \([x,g]\), for \(g\in G\), has finite order? Here \(G\) is a subgroup of the automorphism group, and the fixed-point-free interpretation of the source term ``regular'' requires confirmation.
\subsection{Sources}
\begin{itemize}
\item Supplied local UnsolvedMath registry record, permanent ID \(701138\), exact source code \texttt{catalogue entry}, titled ``Kourovka Notebook Problem 15.50.'' The record labels the problem \texttt{open}; this supplied status has not been independently verified.
\item V. D. Mazurov, Problem 15.50, \emph{The Kourovka Notebook (Unsolved Problems in Group Theory)}, 15th issue (2002), as attributed in the supplied record.
\item The supplied background also identifies the 21st issue edition (2026), arXiv identifier \texttt{1401.0300v48}. The supplied set description names E. I. Khukhro and V. D. Mazurov as editors. These bibliographic details are supplied attribution and have not been independently checked.
\end{itemize}
\end{document}
